\section{Experimental binding data: the first ground-truth label}\label{sec:experimental}
Everything up to this point has been scored \emph{in silico}: the four filters and the
composite rank designs by how AlphaFold3 sees them, with no measurement of whether a design
actually binds. We now have the first experimental readout for the DP3 (known-antibody) set.
This is the non-circular label the scorer has been waiting for (Section~\ref{sec:open}), so
this section sets out exactly what the data is, what it looks like, and --- importantly --- how
far it can be trusted, before any weight is fit to it.

\subsection{How a binding measurement is made}
The DP3 designs were carried into the PepSeq assay. Each library member is displayed as its
constant-length construct (the \num{103}-mer of Section~\ref{sec:tiling}), and the assay reads
out, by sequencing, how much of each member is pulled down by a given monoclonal antibody
relative to a no-antibody control (\texttt{NoAb}) run on the same plate. The working readout is
the one used throughout: $\log_{10}(1+\text{Ab})$ against $\log_{10}(1+\text{NoAb})$, with the
designs cognate to that antibody --- those scaffolding its epitope --- highlighted. A design
that binds sits above the $y=x$ diagonal: more pulldown with the antibody present than without.
We summarize each design by its \emph{cognate log-enrichment},
$\log_{10}(1+\text{Ab}) - \log_{10}(1+\text{NoAb})$ for the antibody raised against its target.

The measurements arrived as two assay runs, and the run parameters are not identical: the six
IM226 antibodies were read at \SI{0.1}{\pico\mol} with antibody-first bead capture, while the
two IM229 antibodies were read at \SI{1}{\pico\mol} (one, 8db4, at a ten-fold higher antibody
concentration) with a different bead order. Absolute counts are therefore not comparable across
runs or even across antibodies, which is the first reason the analysis below stays
\emph{within} an antibody rather than pooling. The full wet-lab protocol is recorded with the
experimental methods; here we treat the intensities as provided.

\subsection{What we have}
Eight monoclonal antibodies, the subset of the original ten that assayed cleanly: 4xwo was
dropped (its synthesized antibody was too low-yield) and 7a3t was dropped (its epitope was too
small for RMSD to discriminate). Their designs remain in the library but have no usable
antibody column. The two runs share the same \num{1000} DP2 library members; \num{403} are our
scaffolded epitope designs (\texttt{category=scaffoldedAbEpitope}) and the rest are control
content (pMHC and published binders, which also carry the only measured $K_d$ values --- none
of our designs do). Of the \num{403}, \num{377} target one of the eight usable antibodies; the
remaining \num{26} target the two dropped ones.

Each design's scaffold sequence matches \texttt{scaffolded\_epitope\_seq} in \texttt{dp2.parquet}
one-to-one (\num{403} of \num{403}), so we can attach its AlphaFold3 metrics directly and ask
how binding relates to them (\texttt{scripts/build\_dp3\_binding\_join.py} $\rightarrow$
\texttt{results/dp3\_binding\_metrics.csv}; figure and numbers from
\texttt{scripts/analyze\_dp3\_binding.py}). Figure~\ref{fig:dp3binding} shows the cognate
designs against the library background for each antibody; Table~\ref{tab:dp3binding} gives the
per-antibody counts and median enrichments.

\begin{figure}[h]\centering
\figorbox{dp3_binding_scatter.png}{1.0}
\caption{DP3 binding, one panel per antibody: $\log_{10}(1+\text{NoAb})$ (x) versus
$\log_{10}(1+\text{Ab})$ (y), all \num{1000} library members in grey, the cognate scaffolded
designs in red, dashed line $y=x$. The two well-sampled antibodies (7ox3, $n=\num{193}$; 5fhx,
$n=\num{149}$) put the bulk of their designs clearly above the diagonal; the small-$n$
antibodies show a handful of points each, mostly above the line. Panels are on independent axes
because absolute counts are not comparable across antibodies/runs.}
\label{fig:dp3binding}
\end{figure}

\begin{table}[h]\centering
\begin{tabular}{llrrr}
\toprule
antibody & run & cognate designs & median log-enrichment & fraction above baseline \\
\midrule
7ox3 & IM226 & 193 & 1.68 & 1.00 \\
5fhx & IM226 & 149 & 0.42 & 0.81 \\
6o9i & IM226 &  12 & 1.44 & 1.00 \\
8pww & IM229 &  11 & 1.85 & 1.00 \\
8jnk & IM226 &   6 & 1.23 & 1.00 \\
8cz8 & IM226 &   3 & 0.86 & 1.00 \\
6xxv & IM226 &   2 & 1.65 & 1.00 \\
8db4 & IM229 &   1 & 0.79 & 1.00 \\
\bottomrule
\end{tabular}
\caption{Per-antibody cognate scaffolded designs. ``Above baseline'' = log-enrichment $>0$
(above the NoAb diagonal). Two antibodies (7ox3, 5fhx) account for \num{342} of the \num{377}
designs; the other six carry \num{35} between them.}
\label{tab:dp3binding}
\end{table}

\subsection{How reliable it is, and how far it goes}
The headline is encouraging: across all eight antibodies the scaffolded designs enrich over the
no-antibody baseline, almost all of them above the diagonal, and for the two well-sampled
antibodies this holds for the great majority of designs. The scaffolding approach produces
binders. That is a real, positive result.

But three properties bound how much this data can teach about \emph{which} design features
predict binding, and they all point the same way --- toward a starting prior, not a fit.

\emph{Everything assayed already passed.} The library was selected for synthesis from the
top of the four-filter ranking, so every one of the \num{377} cognate designs has
\texttt{is\_pass}~$=$~true. The metrics consequently span only their passing range
(\texttt{overall\_rmsd} \numrange{0.35}{1.95}, \texttt{epitope\_chunk\_rmsd} \numrange{0.22}{0.95},
\texttt{mean\_pae} \numrange{3.0}{5.0}, and zero antibody clash by construction). There is no
failing design in the set to contrast against, so the filter itself has no variance here and
cannot be tested for whether it separates binders from non-binders.

\emph{The pooled signal is a between-antibody artifact.} Pooling all \num{377} designs, lower
epitope RMSD tracks higher enrichment (Spearman $\rho=-0.43$), which looks like exactly the
predictor we want. It is mostly an illusion of pooling: different antibodies sit at different
baseline enrichments \emph{and} different metric distributions, so a correlation across the pool
largely reflects which antibody a design belongs to. Within a single antibody the relationship
nearly vanishes (epitope RMSD versus enrichment: $\rho=-0.20$ for 7ox3, $-0.07$ for 5fhx), and
the other metrics are similarly weak and inconsistent in sign.

\emph{The sample is two antibodies deep.} 7ox3 and 5fhx supply \num{342} of the \num{377}
designs; the remaining six antibodies have between one and twelve each. A within-antibody fit is
only meaningfully possible for two epitopes, which is far too narrow to learn weights that would
generalize across targets.

None of this diminishes the positive result --- it constrains its use. DP3 gives a defensible
\emph{starting prior}: the scaffold-versus-linear premise holds, binding is real, and the
cylinder accessibility weight is independently earned by its AUC against the true antibody clash
(Section~\ref{sec:noab}). What DP3 cannot do is fix the composite's weights, because the data has
no failures, a confounded pooled signal, and only two well-sampled epitopes. Closing that gap is
the explicit purpose of the DP4 sampling (Section~\ref{sec:open}): deliberately assay designs
spread across the metric space --- including ones the current filters would reject --- so that
binding can finally be regressed on the metrics with the variance and the breadth this first set
lacks.
